% Master's thesis proposal (English). Build: make proposal → thesis/build/proposal-en.pdf
% Facts and numbers come from the project repository (docs/, src/); sources — SRC-NNN in sources.bib.
\documentclass[a4paper,11pt]{article}
\newcommand{\mainlang}{english}
\newcommand{\docstyle}{article}
% Преамбула: шрифты (латиница, кириллица, арабский с огласовками), поля, библиография.
\usepackage{etoolbox}
\usepackage{fontspec}
\usepackage{polyglossia}
\setdefaultlanguage{\mainlang}
\ifdefstring{\mainlang}{english}{\setotherlanguages{russian,arabic}}{\setotherlanguages{english,arabic}}

% Times New Roman есть в macOS; на другой машине — Liberation Serif или Tempora.
\setmainfont{Times New Roman}
\newfontfamily\cyrillicfont{Times New Roman}
% Арабский — тот же шрифт, что рисует карточки бота (ADR-0024): огласовки и шадда не теряются.
\newfontfamily\arabicfont[Script=Arabic,Path=../src/muallim/presentation/telegram/assets/fonts/,
  Extension=.ttf,UprightFont=*-Regular,BoldFont=*-Bold,Scale=1.25]{ScheherazadeNew}
% \ar{…} — арабская вставка в строке; источник строки — в тексте рядом ([SRC-NNN]).
\newcommand{\ar}[1]{\textarabic{#1}}

\usepackage[left=3cm,right=1.5cm,top=2cm,bottom=2cm]{geometry} % [утвердить] по требованиям вуза
\usepackage{setspace}
\onehalfspacing
\usepackage{indentfirst}

\usepackage{graphicx}
\graphicspath{{../docs/architecture/diagrams/}{../docs/design/screens/}{figures/}}
\usepackage{booktabs,tabularx,longtable}
\usepackage{caption}
\usepackage{csquotes}
\usepackage[hidelinks]{hyperref}

% Библиография — тот же sources.bib, что у всего проекта (ключи SRC-NNN).
\ifdefstring{\mainlang}{english}
  {\usepackage[backend=biber,style=ieee,sorting=none]{biblatex}}
  {\usepackage[backend=biber,style=gost-numeric,sorting=none]{biblatex}}
\addbibresource{../docs/research/sources.bib}

% ID рисунков и таблиц из docs/thesis/figures.md: \label{fig:R-01}, \label{tab:T-01}.

% Данные для титульных листов всех документов. [утвердить] — с руководителем и учебным офисом.
\newcommand{\thesisProgramCode}{09.04.01}                     % [утвердить] код направления
\newcommand{\thesisProgramRu}{Информатика и вычислительная техника} % [утвердить]
\newcommand{\thesisProgramEn}{Computer Science}                % [утвердить]
\newcommand{\thesisSpecialtyRu}{Программная инженерия}
\newcommand{\thesisSpecialtyEn}{Software Engineering}
% Тема — предложение автора (2026-10-18; с ADR-0051 — с адаптацией модели), утверждает руководитель [утвердить].
\newcommand{\thesisTopicRu}{ИИ-тренажёр чтения по букварю «Муаллим ас-Сани»: адаптация открытой модели распознавания фонем к речи русскоязычных начинающих и оценка обратной связи}
\newcommand{\thesisTopicEn}{An AI Reading Tutor for the \emph{Muallim al-Sani} Primer: Adapting an Open Phoneme Recognition Model to Russian-Speaking Beginners and Evaluating Pronunciation Feedback}
\newcommand{\thesisShortRu}{ИИ-тренажёр «Муаллим ас-Сани» и адаптация модели}
\newcommand{\thesisShortEn}{AI Tutor for \emph{Muallim al-Sani} and Model Adaptation}
\newcommand{\thesisAuthorRu}{Исмаилов Абдирахим Ерланович}
\newcommand{\thesisAuthorEn}{Abdirakhim Ismailov}               % латиница — [утвердить] как в паспорте
\newcommand{\thesisAuthorEnFull}{Ismailov Abdirakhim Erlanovich}
\newcommand{\thesisSupervisorRu}{Аскарбекулы Нурсултан}
\newcommand{\thesisSupervisorEn}{Nursultan Askarbekuly}
\newcommand{\thesisUniversityRu}{Университет Иннополис}
\newcommand{\thesisUniversityEn}{Innopolis University}
\newcommand{\thesisCity}{Иннополис}
\newcommand{\thesisCityEn}{Innopolis}
\newcommand{\thesisYear}{2027}

\hypersetup{pdftitle={Thesis Proposal: \thesisShortEn},pdfauthor={\thesisAuthorEn}}
\fancyhead[R]{\small\sffamily\color{inkgray}Thesis proposal · \thesisAuthorEn}
\setcounter{secnumdepth}{2}

\begin{document}
\thispagestyle{plain}
\begin{center}
{\sffamily\small\color{inkgray}\thesisUniversityEn{} · MSc in \thesisSpecialtyEn{} · Master's Thesis Proposal}\\[10pt]
{\sffamily\LARGE\bfseries\color{accent}\thesisTopicEn\par}
\vspace{10pt}
{\color{gold}\rule{0.35\linewidth}{1pt}}\\[8pt]
{\large \thesisAuthorEn}\\[2pt]
{\small Supervisor: \thesisSupervisorEn}\\[2pt]
{\small\color{inkgray} Draft for supervisor review · October 2026 · defence planned for June 2027}
\end{center}

\begin{keybox}[Summary]
Learning to read the Qur'an starts with a primer and a teacher who listens to every syllable.
This thesis builds a Telegram tutor that follows the classical primer \emph{Muallim al-Sani} by Ahmad
Hadi Maksudi page by page, listens to a learner's voice message, detects pronunciation errors at
phoneme level and explains each error in Russian. With the learners' consent the tutor also collects the
first teacher-annotated dataset of primer-level reading by Russian-speaking beginners. On this dataset
the open model Quran Muaalem is \emph{adapted} to the new population in two ways: a calibrated verifier
with a ``not sure'' outcome, and low-rank fine-tuning (LoRA). The work answers four questions: how
closely the off-the-shelf model agrees with teachers and which of its ``learner errors'' are model
errors (RQ-1); how much of the gap to a teacher adaptation on the collected data closes for learners the
model has never seen, without degrading expert recitation (RQ-2, the central question); whether
explanations improve the next attempt and retention (RQ-3); and what the AI-agent harness in which the
system was built and trained prevents and costs (RQ-4). The expected contribution is a measured
adaptation of an open Qur'anic pronunciation model for adult second-language beginners, evaluated on a
frozen, blindly annotated test set; a dataset and a reproducible training pipeline that form the
foundation for doctoral work; and an evidence-based account of building research software with AI
coding agents under mechanical guard-rails.
\end{keybox}

\section{Motivation}

Millions of adults in Russia and Central Asia want to read the Qur'an in Arabic but grew up with
Cyrillic script and without the sounds that Arabic distinguishes. The traditional path is a
primer studied with a teacher. The teacher models a syllable, the learner repeats it, and the teacher
corrects every slip on the spot. This oral transmission, \emph{talaqqi}, works well, but it does not
scale: a teacher can listen to one learner at a time, and many learners live far from a mosque school
or study after work.

One primer has dominated this region for more than a century. Ahmad Hadi Maksudi (1868--1941)
wrote \emph{Muallim al-Sani} (``The Second Teacher'') as a bridge from knowing the letters to reading
the Qur'an correctly. It follows the sound-based method of the Jadid reformers: a letter is learned
as a sound, not as a name, and syllables grow into words in small steps
\cite{SRC-015,SRC-016,SRC-017}. The book is still reprinted and officially distributed by Muslim
boards in Tatarstan and Kazakhstan \cite{SRC-020,SRC-040}.

Speech technology has meanwhile reached the point where a machine can listen at the level of
phonemes. Quran Muaalem, an open model released in 2025, recognises phonemes together with ten
articulatory attributes (\emph{sifat}) and reports a phoneme error rate of 0.54\,\% on expert
recitations \cite{SRC-002}. Meta-analyses show that automatic pronunciation feedback has a medium
positive effect on learners ($g=0.69$), largest for individual sounds and explicit corrective
feedback \cite{SRC-025}, which is precisely the task of a primer. The opportunity is clear: a tutor
that gives every beginner a patient listener between lessons with a human teacher.

The obstacle is equally clear. Muaalem was trained only on male expert reciters. On the
qdat\_bench set with real errors its phoneme error rate rises to 5.78\,\%, and the F1 score for the
\emph{ikhfa} rule falls to 0.45 \cite{SRC-002}. The best system in the IQRA 2026 challenge reaches a
phoneme-level F1 of 0.72, and the organisers note that no second-language learners are present in the
data \cite{SRC-031,SRC-086}. Automatic systems also produce false alarms on accented speech, which is
why reviews recommend keeping a teacher in the loop \cite{SRC-026}. Nobody has measured how such a
model behaves on a Russian-speaking adult who reads isolated syllables slowly, for the first time,
and nobody has tried to adapt it to such a learner. Adaptation is feasible with little data:
fine-tuning self-supervised speech encoders on learner speech improves mispronunciation detection
\cite{SRC-116}, low-rank adaptation changes only a small fraction of the weights \cite{SRC-111}, and
calibration with selective prediction lets a model say ``not sure'' instead of issuing a false verdict
\cite{SRC-112,SRC-117}.

\section{Research gap}

A review of literature, applications and datasets carried out for this proposal
\cite{SRC-038} found five gaps that this thesis can address within a master's timeline.

\begin{enumerate}
  \item \textbf{No evaluation on primer-level beginners.} Benchmarks for Qur'anic pronunciation
    (QuranMB, Iqra'Eval, IQRA 2026) contain fluent recitation of verses
    \cite{SRC-004,SRC-043,SRC-031}. The closest dataset of non-Arab learners, RetaSy, contains
    no Russian or Kazakh speakers \cite{SRC-032}.
  \item \textbf{No studies of Russian- and Kazakh-speaking learners' errors.} The typical
    difficulties (pharyngeals, emphatics, interdentals, vowel length, gemination) are described in
    methodological sources and in studies of other first languages \cite{SRC-023,SRC-034,SRC-035},
    but not measured for this population.
  \item \textbf{No adaptation of models to second-language beginners.} Open Qur'anic models are trained
    on experts \cite{SRC-002}; we found no work that fine-tunes such a model on non-Arab beginners and
    honestly measures the gain on unseen learners, while small-data ML studies frequently leak
    information between training and evaluation \cite{SRC-122}.
  \item \textbf{No digital adaptation of the primer with automatic checking,} and no controlled
    trials of primer-based methods for reading the Qur'an in general \cite{SRC-038}. Telegram bots for
    pronunciation have been studied only as channels, not as automatic tutors \cite{SRC-092}.
  \item \textbf{Weak evidence on AI-agent software engineering.} Agents are known to game their own
    tests \cite{SRC-027,SRC-082}; repository context files barely change correctness
    \cite{SRC-074,SRC-077}; practitioners describe ``harness engineering'' \cite{SRC-029,SRC-028},
    but few studies report how a guarded harness behaves over a whole research project
    \cite{SRC-073}.
\end{enumerate}

\section{Aim, objectives and research questions}

\textbf{Aim.} To build a reading tutor that faithfully follows \emph{Muallim al-Sani}, to collect
with it an annotated dataset of Russian-speaking beginners' reading, and to adapt an open phoneme-level
model on this dataset so that it takes the role of the listening teacher more reliably --- measuring how
much of the gap to a teacher the adaptation closes and where the gap remains.

\textbf{Objectives.}
\begin{enumerate}[label=O\arabic*.]
  \item Turn the primer into a digital curriculum in the book's own order, with every exercise traceable
    to a page and line of a concrete edition \cite{SRC-040}.
  \item Build a tutor that assesses phoneme-level reading, explains errors in Russian for a learner who
    starts from zero, and records a full decision trace for each attempt.
  \item Collect a dataset: consent (including consent to model training) and pseudonyms, a closed pilot
    with 20--40 learners over 4--6 weeks, phoneme-level annotation by two teachers, blind annotation
    of test learners, and a datasheet \cite{SRC-123}.
  \item Build a reproducible training and evaluation pipeline: a learner-disjoint split with a frozen
    test set, the calibrated verifier B1, LoRA fine-tuning B2 with augmentation from the primer's
    minimal pairs, and a protected evaluator.
  \item Run a within-subject feedback experiment and a delayed post-test.
  \item Answer RQ-1 to RQ-3 with effect sizes and confidence intervals following the preregistered plan.
  \item Document the AI-agent harness in which the system was built and trained, and evaluate it as a
    case study (RQ-4).
\end{enumerate}

\begin{rqbox}
\textbf{RQ-1 (baseline).} How closely does the off-the-shelf assessment B0 --- Muaalem, rule-based
alignment and grading thresholds --- agree with teachers on beginners reading \emph{Muallim al-Sani}, and
which of the errors it reports belong to the learner rather than to the model?\\[4pt]
\textbf{RQ-2 (adaptation, central).} How much does adapting the model on the collected dataset --- a
calibrated verifier with a ``not sure'' outcome (B1) and low-rank fine-tuning (B2) --- reduce false
rejections and bring the assessment closer to teachers on learners unseen in training, without
degrading recognition of expert recitation?\\[4pt]
\textbf{RQ-3 (feedback).} Do explicit explanations in Russian improve the next attempt and delayed
retention compared with a bare correct/incorrect signal?\\[4pt]
\textbf{RQ-4 (method).} Which known failure modes of AI-agent development and ML experimentation does a
mechanically enforced harness prevent or expose, and how much human intervention does it require?
\end{rqbox}

The hypotheses below are recorded in a preregistration that will be frozen with a SHA-256
fingerprint before data collection begins \cite{SRC-067}. Everything else is reported as exploratory.

\begin{table}[ht]
\centering\small
\caption{Confirmatory hypotheses (preregistration v0.3, to be approved with the supervisor).}
\label{tab:hyp-en}
\begin{tabularx}{\linewidth}{@{}lL@{}}
\toprule
ID & Hypothesis \\
\midrule
H1.1 & Precision and recall of B0 error detection on beginners are lower than reported on the model's
       training domain (Qur'anic recitation). \\
H1.2 & For B0, false rejections (FRR) occur more often than false acceptances (FAR). \\
H1.3 & B0--teacher agreement ($\kappa$) is lower than teacher--teacher agreement. \\
H1.4 & Russian speakers make more confirmed errors in emphasis (\emph{tafkhim}) than in place of
       articulation. \\
H2.1 & On test learners, B1 has a lower FRR than B0 while its recall is no lower than B0's minus 0.05. \\
H2.2 & On test learners, B2 has a higher phoneme-level F1 than B0. \\
H2.3 & B2 is not worse than B0 on expert recitation: its phoneme error rate rises by at most one
       percentage point. \\
H3.1 & After a full explanation, the score gain on the next attempt is larger than after a
       correct/incorrect signal. \\
H3.2 & The difference in H3.1 persists on a delayed test after 1--2 weeks. \\
\bottomrule
\end{tabularx}
\end{table}

\paragraph{Scope} The thesis covers the whole primer as curriculum (pages 10--85), but the
evaluation concentrates on the letter and vowel lessons (pages 10--24), where beginners start and where
errors are most frequent. Tajweed rules such as \emph{ikhfa} and \emph{idgham}, other first languages,
active learning and per-learner personalisation of the model are left for doctoral work. The model is
adapted, not trained from scratch: the pilot yields orders of magnitude less data than training from
scratch would need. The tutor complements a teacher; it does not claim to replace one.

\section{Research approach}

The thesis follows design science research \cite{SRC-100,SRC-101}: it builds artefacts that solve a
real problem and evaluates them with methods suited to each claim. Three artefacts are involved
(Figure~\ref{fig:design-en}). Artefact~A is the reading tutor together with its dataset, evaluated by a
measurement study and a pilot experiment. Artefact~C is the adapted model and its training pipeline,
evaluated by a controlled experiment on a frozen test set. Artefact~B is the development harness,
evaluated as an embedded case study following Runeson and Höst \cite{SRC-102}. The ACM SIGSOFT Empirical Standards
serve as a checklist for each study type \cite{SRC-103}.

\begin{figure}[ht]
\centering
\resizebox{\linewidth}{!}{% Дизайн исследования: три артефакта, четыре RQ, данные и тип исследования.
% Источник: docs/research/metrics-gqm.md (G1–G7), preregistration.md (v0.3), roadmap 1.4–1.8, ADR-0051.
\begin{tikzpicture}[node distance=4mm and 6mm]
  \node[hot,text width=4.1cm] (a1) {\textbf{\LANG{Artefact A}{Артефакт A}}: \LANG{reading tutor}{тренажёр чтения}\\[-1pt]{\scriptsize\LANG{book-ordered curriculum, B0 assessment, feedback, dataset}{программа по книге, оценка B0, объяснения, набор данных}}};
  \node[hot,right=5mm of a1,text width=4.1cm] (a3) {\textbf{\LANG{Artefact C}{Артефакт C}}: \LANG{adapted model}{адаптированная модель}\\[-1pt]{\scriptsize\LANG{B1 calibrated verifier, B2 LoRA fine-tuning}{B1 калиброванный верификатор, B2 дообучение LoRA}}};
  \node[goldbox,right=5mm of a3,text width=4.1cm] (a2) {\textbf{\LANG{Artefact B}{Артефакт B}}: \LANG{agent harness}{обвязка агента}\\[-1pt]{\scriptsize\LANG{hooks, frozen tests and test set, reviewers, journals}{хуки, замороженные тесты и тест, ревьюеры, журналы}}};

  \node[rqbox-node,minimum height=1.25cm,below=7mm of a1,xshift=-2.0cm,text width=2.9cm] (rq1) {\textbf{RQ-1}\\[-1pt]{\scriptsize\LANG{baseline vs.\ teachers; error profile}{B0 против учителей; профиль ошибок}}};
  \node[rqbox-node,minimum height=1.25cm,right=of rq1,text width=2.9cm] (rq2) {\textbf{RQ-2}\\[-1pt]{\scriptsize\LANG{does adaptation close the gap?}{закрывает ли адаптация разрыв?}}};
  \node[rqbox-node,minimum height=1.25cm,right=of rq2,text width=2.9cm] (rq3) {\textbf{RQ-3}\\[-1pt]{\scriptsize\LANG{effect of explanations}{эффект объяснений}}};
  \node[rqbox-node,minimum height=1.25cm,right=of rq3,text width=2.9cm] (rq4) {\textbf{RQ-4}\\[-1pt]{\scriptsize\LANG{what the harness prevents, at what cost}{что предотвращает обвязка и какой ценой}}};

  \node[soft,minimum height=1.75cm,below=of rq1,text width=2.9cm] (d1) {\scriptsize\LANG{phoneme-level annotation; P/R/F1, FAR/FRR, $\kappa$; confirmed confusions}{разметка по фонемам; P/R/F1, FAR/FRR, $\kappa$; подтверждённые путаницы}};
  \node[soft,minimum height=1.75cm,below=of rq2,text width=2.9cm] (d2) {\scriptsize\LANG{learner-disjoint frozen test, blind gold labels; B0 vs.\ B1 vs.\ B2}{замороженный тест по ученикам, слепое «золото»; B0, B1, B2}};
  \node[soft,minimum height=1.75cm,below=of rq3,text width=2.9cm] (d3) {\scriptsize\LANG{within-subject: full vs.\ minimal feedback; delayed test}{внутрисубъектно: полное vs.\ краткое; отложенный тест}};
  \node[soft,minimum height=1.75cm,below=of rq4,text width=2.9cm] (d4) {\scriptsize\LANG{embedded case study: hook logs, sessions, commits}{встроенное кейс-стади: журналы хуков, сессии, коммиты}};

  \node[lab,below=1mm of d1] {\LANG{measurement study}{измерительное}};
  \node[lab,below=1mm of d2] {\LANG{controlled ML experiment}{контролируемый ML-эксперимент}};
  \node[lab,below=1mm of d3] {\LANG{pilot experiment}{пилотный эксперимент}};
  \node[lab,below=1mm of d4] {\LANG{case study}{кейс-стади}};

  \draw[arr] (a1.south) -- (rq1.north);
  \draw[arr] (a1.south) -- (rq3.north);
  \draw[arr] (a1.south) -- (rq2.north);
  \draw[arr] (a3.south) -- (rq2.north);
  \draw[arr] (a2.south) -- (rq4.north);
  \foreach \r/\d in {rq1/d1,rq2/d2,rq3/d3,rq4/d4}{\draw[arr] (\r) -- (\d);}
\end{tikzpicture}
}
\caption{Research design: three artefacts, four research questions, and the evidence for each.}
\label{fig:design-en}
\end{figure}

\subsection{Artefact A: the reading tutor}

The tutor already exists as a working system (version 0.6.0). The curriculum contains the entire
book --- 125 lessons and 800 cards in nine parts --- in the book's order, from the hamza and the group
\ar{ر ز م} on page~10 to the short suras at the end. Each card stores the page and line of the 2010
Almetyevsk edition, and the 139 Qur'anic examples were matched to the Medina mushaf text
\cite{SRC-040,SRC-041}.

A learner opens a card, may listen to a reference reading, and sends a voice message. The assessment
pipeline (Figure~\ref{fig:pipe-en}) converts the audio, recognises phonemes with Muaalem, aligns them
with the canonical phonemes of the line for both pause and continuation readings, classifies
differences into ten error kinds tied to a syllable, and applies grading thresholds. The learner sees
the line with the error highlighted, a short ``how it should sound $\to$ how you said it'' comparison, an
articulation hint in Russian, and a next step (``listen, say the syllable slowly, read the line
again''). This design draws on findings that explicit, focused feedback works better than recasts and
long lists \cite{SRC-053,SRC-054,SRC-055,SRC-056}.

\begin{figure}[ht]
\centering
\resizebox{\linewidth}{!}{% Конвейер оценки одной попытки: от голосового до объяснения и записи для исследования.
% Код: application/use_cases/recitation.py, infrastructure/asr/muaalem.py, domain/assessment.py,
% domain/policy.py, application/feedback.py (ADR-0007, ADR-0014, ADR-0019).
\begin{tikzpicture}[node distance=5mm and 7mm]
  \node[goldbox,text width=2.7cm] (voice) {\LANG{Voice message}{Голосовое}\\[-1pt]{\scriptsize OGG/Opus, $\le$60\,s}};
  \node[box,right=of voice,text width=2.6cm] (ff) {ffmpeg\\[-1pt]{\scriptsize\LANG{16 kHz mono}{16 кГц, моно}}};
  \node[hot,right=of ff,text width=3.0cm] (asr) {Quran Muaalem\\[-1pt]{\scriptsize\LANG{phonemes + 10 \emph{sifat}}{фонемы + 10 сифат}}};

  \node[goldbox,below=11mm of voice,text width=2.7cm] (card) {\LANG{Book line}{Строка книги}\\[-1pt]{\scriptsize\LANG{page + line, SRC-040}{страница + строка, SRC-040}}};
  \node[box,right=of card,text width=2.6cm] (phon) {\LANG{Phonetiser}{Фонетизатор}\\[-1pt]{\scriptsize\LANG{\emph{wasl} and \emph{waqf}}{васл и вакф}}};
  \node[box,right=of phon,text width=3.0cm] (align) {\LANG{Alignment}{Выравнивание}\\[-1pt]{\scriptsize\LANG{per reference variant}{с каждым эталоном}}};

  \node[box,right=of asr,text width=2.9cm,yshift=-8mm] (err) {\LANG{Typed errors}{Типизированные ошибки}\\[-1pt]{\scriptsize\LANG{10 kinds, syllable}{10 типов, слог}}};
  \node[box,below=5mm of err,text width=2.9cm] (pol) {\LANG{Grading policy}{Политика оценки}\\[-1pt]{\scriptsize\LANG{score, verdict}{балл, вердикт}}};
  \node[human,below=5mm of pol,text width=2.9cm] (fb) {\LANG{Feedback in Russian}{Объяснение по-русски}\\[-1pt]{\scriptsize\LANG{what / how / next}{что / как / что дальше}}};

  \node[store,below=9mm of align,text width=3.1cm,minimum height=1.2cm] (db) {\LANG{Attempt record}{Запись попытки}\\[-1pt]{\scriptsize\LANG{full decision trace}{полный след решения}}};

  \draw[arr] (voice) -- (ff);
  \draw[arr] (ff) -- (asr);
  \draw[arr] (card) -- (phon);
  \draw[arr] (phon) -- (align);
  \draw[arr] (asr.south) -- node[lab,right]{\LANG{heard}{услышано}} (align.north);
  \draw[arr] (align.east) -- ++(4mm,0) |- (err.west);
  \draw[arr] (err) -- (pol);
  \draw[arr] (pol) -- (fb);
  \draw[darr] (align) -- (db);
  \node[lab,below=1mm of db] {\LANG{audio kept only with consent (SHA-256)}{аудио — только с согласия (SHA-256)}};
\end{tikzpicture}
}
\caption{Assessment of one attempt, from voice message to explanation and research record.}
\label{fig:pipe-en}
\end{figure}

Every attempt keeps the reference variants, the recognised phonemes, the raw model output, thresholds,
code, curriculum and model versions, so that any verdict can be recomputed later with different
thresholds or a different model. Audio is stored only with explicit, revocable consent; learners are
identified by pseudonyms; and a learner can delete all their data with one button.

\subsection{Artefact C: the adapted model}

Three systems are compared on the same recordings. \textbf{B0} is Muaalem as is, with rules and
thresholds; this is how the tutor works today. \textbf{B1} is a verifier trained on top of B0's outputs:
for every phoneme it receives the model's confidence, the probabilities of the articulatory attributes,
the Goodness of Pronunciation score \cite{SRC-115}, duration and material type, and estimates the
probability of a real error. Its probabilities are calibrated \cite{SRC-112}, and its threshold is chosen
by conformal risk control so that the false-rejection rate stays below a chosen $\varepsilon$
\cite{SRC-113,SRC-114}; attempts without enough confidence receive a ``not sure'' outcome
\cite{SRC-117} --- the learner is asked to read again and the attempt goes to a teacher. The verifier
needs only hundreds of examples and trains on a CPU. \textbf{B2} is Muaalem fine-tuned with LoRA
\cite{SRC-111} on teacher-corrected phonemes, as self-supervised speech encoders have been fine-tuned for
learner mispronunciation detection before \cite{SRC-116}. To obtain more errors for training, the book
itself is used: a correct reading of a card with \ar{ك} is a ready-made error example for the paired
card with \ar{ق}, and likewise for every minimal pair. This is a form of phoneme-substitution
augmentation \cite{SRC-119}, complemented by speed perturbation \cite{SRC-120}. The fine-tuned model is
also checked on expert recitation, so that gains on beginners are not bought by forgetting
\cite{SRC-118}.

At the end of the pilot, B1 and B2 run inside the tutor \emph{in shadow mode}: the learner still sees
B0's verdict, while the adapted systems' verdicts are recorded alongside it. This keeps the RQ-3
experiment free of a model change and gives an additional post-training check on new attempts.

\subsection{Artefact B: the development harness}

The system is built by one author with an AI coding agent (Claude Code). Rules that matter are
enforced by the environment, not by the agent's goodwill \cite{SRC-029,SRC-076}. Hooks block edits
to tests during the implementation phase, refuse reads of personal data, ask for human approval in
human-owned zones (curriculum, texts shown to learners, thresholds, the RQ-1 evaluator), check that a
new library really exists on PyPI, and stop the agent from finishing while architecture contracts, type
checks or tests are red. Reviewer subagents without write access check code, architecture and whether
the change matches what the author asked for. The same mechanisms protect training experiments: the
list of test learners and the evaluator are frozen with a SHA-256 fingerprint, and a hook forbids the
agent to read test labels, following the rules formulated by the supervisor's group for research with
agents \cite{SRC-067}. Each block and each approval request is written to a journal. This journal, together with the session and prompt logs, is the data for RQ-4.

\subsection{Studies}

\paragraph{Study 1 --- agreement (RQ-1)} Two teachers independently annotate every pilot attempt at
phoneme level, seeing the tutor's answer and correcting it where it is wrong. Following the Iqra'Eval
and L2-ARCTIC protocols \cite{SRC-063,SRC-064}, each reference phoneme is classified as true
acceptance, true rejection, false acceptance or false rejection, giving precision, recall, F1 and
diagnosis accuracy. At verdict level we compute FAR and FRR, and Cohen's $\kappa$ \cite{SRC-104}
between the tutor and each teacher, compared with $\kappa$ between the two teachers as a ceiling.
Agreement between annotators per phoneme is measured with Krippendorff's $\alpha$ \cite{SRC-065}.
Confidence intervals come from a bootstrap that resamples learners, because attempts by one learner are
not independent. Self-corrections and unusable recordings are reported separately.

The error profile belongs to the same study. Confusion matrices ``expected $\to$ heard'' are built
twice: for all errors reported by the tutor, and only for errors confirmed by annotators. The difference
shows how much of the apparent learner profile is model noise. The primer's own minimal pairs
(\ar{ق}/\ar{ك}, \ar{ص}/\ar{س}, \ar{ط}/\ar{ت}, \ar{ح}/\ar{خ}, \ar{ع}/\ar{غ} and others) serve as a
pre-defined lens. A short entry questionnaire records the learner's first language and prior reading
experience as covariates.

\paragraph{Study 2 --- model adaptation (RQ-2)} Before the first training run, about 30\,\% of learners,
stratified by first language and activity, are assigned to the test set, and the list is fingerprinted.
No recording of a test learner is used for training, calibration or threshold selection: splitting by
attempt would leak the speaker's voice \cite{SRC-122}. Both teachers annotate the test learners'
attempts \emph{blind}, without seeing the tutor's answer, and reconcile disagreements; the resulting gold
labels do not tilt the comparison towards B0. Hyperparameters are chosen only by learner-grouped
cross-validation inside the training part, since tuning on the evaluation data inflates estimates
\cite{SRC-121}. The test set is opened once, after all systems are trained and frozen. B0, B1 and B2 are
compared with the metrics of Study~1 using a paired cluster bootstrap over learners. The summary measure
is the \emph{share of the gap closed}: how far a system's $\kappa$ with the gold labels rises from B0
towards the teacher--teacher ceiling. Calibration error, the verifier's risk--coverage curve and, as an
exploratory analysis, a learning curve --- how B2 improves with the number of training learners --- are
also reported. The learning curve answers a doctoral question directly: how much data a new learner
population needs.

\paragraph{Study 3 --- feedback experiment (RQ-3)} Because 20--40 learners are too few for two
independent groups (a medium effect needs about 45 per group), the design is within-subject, after
Neri et al.\ \cite{SRC-054}. Each learner receives a full explanation on a random half of the cards
and only a correct/incorrect signal on the other half. The primary outcome is the score change from a
failed attempt to the next attempt on the same card; the secondary outcome is the score on a delayed
test after 1--2 weeks. A mixed model with learner and card as random effects estimates the effect
\cite{SRC-106}. Holm correction is applied across the nine confirmatory hypotheses.

\paragraph{Study 4 --- harness case study (RQ-4)} Data sources are the hook journal (blocks and
approval requests by reason), session summaries, prompt logs, commits and architecture decision records.
Measures include blocks by type, human interventions per session, iterations until green checks,
defects that the checks missed, and attempts by the agent to access the test set. Self-reported productivity is not used, because it diverges from
measured productivity \cite{SRC-030}.

\section{Expected contributions}

\begin{enumerate}
  \item \textbf{Empirical:} the first teacher-referenced measurement of an open Qur'anic phoneme model
    on primer-level reading by second-language beginners, including where it fails and how often it
    wrongly rejects correct reading.
  \item \textbf{Model and method:} small-data adaptation of an open model to a new population --- a
    calibrated verifier with a guarantee on the false-rejection rate and LoRA fine-tuning with
    augmentation from the primer's minimal pairs --- honestly evaluated on unseen learners and checked
    for forgetting.
  \item \textbf{Data:} a consented set of beginners' readings with phoneme-level labels, a blind gold
    test set and a datasheet, released together with the fine-tuned weights if ethics approval and
    consent allow \cite{SRC-097,SRC-123}.
  \item \textbf{Software-engineering method:} an evidence-based description of a guarded AI-agent
    harness for research software and ML experiments, with its measured benefits, costs and limits.
  \item \textbf{Artefact:} an open-architecture tutor where model, database and messenger can be
    replaced without touching the domain logic --- a ready test bed for comparing models.
\end{enumerate}

\paragraph{Foundation for doctoral work} The master's thesis leaves the doctoral project not only
questions but also instruments: a frozen test set that remains the test set, a training pipeline, a
learning curve and shadow mode. The main doctoral line follows from them --- human-in-the-loop
adaptation: which attempts to send to a teacher for labelling (active learning), how to use ``the bot
was wrong'' reports as weak labels, how to personalise the model per learner and transfer it to other
first languages.

\section{Ethics, data protection and religious sensitivity}

Voice recordings combined with a religious topic make the data sensitive. The design therefore
separates consent to the service from consent to voice recording, never stores audio without
consent, uses pseudonyms in every report and export, restricts research participation to adults,
and gives every learner a one-button deletion \cite{SRC-033,SRC-036,SRC-037}. The open questions are
institutional: who acts as data operator, whether the university ethics committee must review the
study, and where the server is located. They will be settled with the supervisor before any learner
outside the closed circle is invited. The tutor uses only the Arabic text of the mushaf and its own
neutral explanations; it does not reproduce the book's Russian prose, which is protected by copyright.
Consent to voice recording explicitly covers model training and the release of weights; no voice prints
are computed. Muaalem and its base model w2v-BERT~2.0 are released under the MIT licence \cite{SRC-003},
so fine-tuning and releasing weights are permitted.

\section{Work plan}

Figure~\ref{fig:gantt-en} shows the plan. Preparation (proposal, ethics, proofreading of the curriculum
against the scan, reference recordings, recruiting annotators) runs until December 2026; in November
and December the training pipeline is debugged on open data, so that in April it only has to be run on
pilot data. Data collection runs from January to March 2027, and the list of test learners is frozen in
February. Training B1 and B2 and the single test run take place in April; analysis and writing follow,
with a pre-defence in May and the defence in June 2027.

\begin{figure}[ht]
\centering
\resizebox{\linewidth}{!}{% План работ: октябрь 2026 — июнь 2027 (roadmap, этап 1). Месяц = 2 слота (половины месяца).
\begin{ganttchart}[
    hgrid,vgrid={*1{draw=black!8}},
    x unit=0.68cm,y unit chart=0.52cm,y unit title=0.55cm,
    title/.append style={fill=accent!10,draw=accent!40},
    title label font=\sffamily\scriptsize,
    bar/.append style={fill=accent!70,draw=accent},
    bar height=0.55,
    group/.append style={fill=gold,draw=gold},
    group height=0.25,
    milestone/.append style={fill=lesson,draw=lesson},
    bar label font=\sffamily\footnotesize,
    group label font=\sffamily\footnotesize\bfseries,
    milestone label font=\sffamily\footnotesize\itshape,
  ]{1}{18}
  \gantttitle{\LANG{Oct}{Окт}}{2}\gantttitle{\LANG{Nov}{Ноя}}{2}\gantttitle{\LANG{Dec}{Дек}}{2}
  \gantttitle{\LANG{Jan}{Янв}}{2}\gantttitle{\LANG{Feb}{Фев}}{2}\gantttitle{\LANG{Mar}{Мар}}{2}
  \gantttitle{\LANG{Apr}{Апр}}{2}\gantttitle{\LANG{May}{Май}}{2}\gantttitle{\LANG{Jun}{Июн}}{2}\\
  \ganttgroup{\LANG{Preparation}{Подготовка}}{1}{6}\\
  \ganttbar{\LANG{Proposal, preregistration freeze}{Предложение, заморозка предрегистрации}}{1}{3}\\
  \ganttbar{\LANG{Ethics, consent, legal debt (T1)}{Этика, согласие, правовой долг (T1)}}{2}{5}\\
  \ganttbar{\LANG{Book proofreading, reference audio}{Сверка книги, эталонные записи}}{1}{5}\\
  \ganttbar{\LANG{Annotators, pilot readiness}{Разметчики, готовность пилота}}{4}{6}\\
  \ganttgroup{\LANG{Data collection}{Сбор данных}}{7}{12}\\
  \ganttbar{\LANG{Pilot, 20--40 learners, 6 weeks}{Пилот, 20--40 учеников, 6 недель}}{7}{9}\\
  \ganttbar{\LANG{Delayed post-test (RQ-3)}{Отложенный тест (RQ-3)}}{10}{10}\\
  \ganttbar{\LANG{Double annotation (RQ-1)}{Двойная разметка (RQ-1)}}{8}{12}\\
  \ganttbar{\LANG{Agent-harness telemetry (RQ-4)}{Журналы обвязки агента (RQ-4)}}{1}{12}\\
  \ganttgroup{\LANG{Analysis and writing}{Анализ и текст}}{11}{18}\\
  \ganttbar{\LANG{Analysis RQ-1--RQ-4}{Анализ RQ-1--RQ-4}}{11}{14}\\
  \ganttbar{\LANG{Chapters 4--6, abstract}{Главы 4--6, аннотация}}{12}{16}\\
  \ganttmilestone{\LANG{Pre-defence}{Предзащита}}{16}\\
  \ganttbar{\LANG{Revision, plagiarism check}{Доработка, проверка заимствований}}{16}{17}\\
  \ganttmilestone{\LANG{Defence}{Защита}}{18}
\end{ganttchart}
}
\caption{Work plan, October 2026 -- June 2027 (half-month resolution).}
\label{fig:gantt-en}
\end{figure}

\section{Risks}

\begin{table}[ht]
\centering\small
\caption{Main risks and responses (from the project risk register).}
\begin{tabularx}{\linewidth}{@{}P{4.1cm}L@{}}
\toprule
Risk & Response \\
\midrule
Too few pilot learners & Recruit through teachers and mosques early; within-subject design needs
fewer people; report effect sizes with intervals rather than significance alone. \\
The model often errs on beginners & This is itself an RQ-1 result if measured honestly; a
``the bot was wrong'' button, raw model outputs in every attempt, and a planned ``not sure'' outcome. \\
No second annotator & Agree with two teachers before the pilot; fallback: the author plus one teacher,
both labelling every attempt. \\
Ethics and data-protection blockers & Raise operator and ethics questions with the supervisor now;
close all pilot-blocking items in the technical-debt register first. \\
Too little data for fine-tuning & B1 works with hundreds of examples; minimal-pair augmentation; if
fewer than about 1\,500 annotated attempts exist by 15 April, B2 is reported as a feasibility check and
RQ-2 is answered by B1 (fixed in the preregistration). \\
Leakage between training and test & Learner-disjoint split; test list and evaluator frozen by
fingerprint; a hook forbids the agent to read test labels; the test set is opened once. \\
No GPU & University cluster or cloud; LoRA needs a single card; B1 trains on a CPU. \\
Errors in the Arabic text of the curriculum & Human proofreading against the scan; every attempt
stores the curriculum version, so affected attempts can be excluded. \\
Scope larger than the timeline & Prioritise RQ-1 and RQ-2 on pages 10--24; MoSCoW priorities; review at each
supervisor meeting. \\
\bottomrule
\end{tabularx}
\end{table}

\section{Current status and preliminary evidence}

The system runs end to end with the real model. In a first sanity check on the teaching recitation of
Mahmoud Khalil al-Husary \cite{SRC-042}, correct verses scored 99--100 and a deliberately mismatched
verse scored 40--67, so the pipeline separates right from wrong reading on expert audio. On nine early
attempts the median processing time was 1.7\,s (90th percentile 5.0\,s), within the 10-second target,
although this excludes download and rendering. Properties of the assessment algorithm were checked on
5\,000 random phoneme strings per property without a violation \cite{SRC-096}. The code base has about
9\,000 lines of Python in 92 modules, 736 passing automated tests, 50 architecture decision records and
47 functional requirements traced to tests. Every attempt already stores the raw model output with
per-phoneme confidence and the model version --- exactly what training B1 and comparing systems
requires. None of these numbers answers a research question yet;
they show that the instrument is ready.

\section{Proposed thesis outline}

\begin{enumerate}[leftmargin=2em]
  \item \textbf{Introduction} --- context, problem, aim, research questions, contributions.
  \item \textbf{Literature Review} --- the primer and the sound-based method; computer-assisted
    pronunciation training; Qur'anic mispronunciation detection; adapting and calibrating speech
    models; AI-agent software engineering.
  \item \textbf{Methodology} --- design science; requirements and quality scenarios; GQM metrics;
    study designs; data splitting and test-set protection; preregistration; threats to validity.
  \item \textbf{Implementation} --- curriculum from the book; architecture; assessment and feedback;
    research data and annotation; the B1 and B2 training pipeline; privacy; the agent harness.
  \item \textbf{Evaluation and Discussion} --- RQ-1 to RQ-4, limitations.
  \item \textbf{Conclusion} --- answers, contributions, doctoral directions.
\end{enumerate}

\printbibliography[heading=bibintoc,title={References}]
\end{document}
